% ===========================================================================
%  GENERAL CONCLUSION
% ===========================================================================
\chapter*{General Conclusion}
\addcontentsline{toc}{chapter}{General Conclusion}
\markboth{General Conclusion}{General Conclusion}


\section*{Summary of Findings}

This thesis investigated whether artificial intelligence and machine learning methods can improve the timeliness and accuracy of France's Balance of Payments estimates, and whether such improvements matter for macroeconomic policy.  Three self-contained chapters provided complementary answers.

\textbf{Chapter~1} conducted a systematic horse race between eight models spanning autoregressive benchmarks, econometric approaches (DFM, Bridge), regularised linear ML (Ridge), nonlinear ML (Gradient Boosting, XGBoost), and deep learning (LSTM).  In a pseudo-real-time expanding-window evaluation on quarterly data (2000--2022), \textbf{Ridge regression achieved the largest improvement: a~26.3\% RMSE reduction over AR(1)} ($p = 0.002$).  Traditional econometric models delivered moderate gains (DFM~$-6.6\%$, Bridge~$-13.8\%$), nonlinear ML improved on the benchmark but was dominated by Ridge, and LSTM failed entirely ($+1.0\%$).  The key takeaway: in small-sample macroeconomic settings, the bias--variance tradeoff favours regularisation over complexity.

\textbf{Chapter~2} conducted a structured ablation study to quantify the incremental value of four categories of alternative data: trade-activity proxies, financial-flow indicators, text/sentiment data, and satellite/geospatial proxies.  The \textbf{full model achieved a statistically significant~6.5\% RMSE reduction} over baseline Ridge (DM $p = 0.020$), with \textbf{text/sentiment data providing the largest SHAP contribution (16.8\%)}.  No individual category was significant in isolation---but their combination, optimally aggregated by Ridge's $\ell_2$ penalty, was.  All eight data sources come from real public APIs, ensuring full reproducibility.

\textbf{Chapter~3} bridged accuracy and policy relevance through three exercises.  Crisis evaluation showed Ridge maintained~18--20\% RMSE reductions during COVID and the energy shock.  Cross-country tests confirmed generalisability: \textbf{16--20\% improvements for Germany, Italy, and Spain}.  A BoP-augmented Taylor rule estimated $\hat{\delta} = -0.481$ ($p = 0.003$), implying that \textbf{more accurate BoP nowcasts would have shifted implied policy rates by up to~159.7 basis points} at peak episodes.


\section*{Contributions}

The thesis makes five contributions to the literature:

\begin{enumerate}
  \item \textbf{First systematic ML-based BoP nowcasting study.}  While GDP nowcasting has been extensively studied \citep{giannone2008nowcasting, bok2018macroeconomic}, this is---to our knowledge---the first paper to conduct a rigorous model comparison for Balance of Payments estimation.

  \item \textbf{Quantification of alternative data value.}  The ablation design with nested models (M0--M5) provides the first structured evidence on which alternative data categories matter for external sector prediction, with statistical testing and interpretability via SHAP.

  \item \textbf{Policy counterfactual.}  The Taylor rule exercise provides a concrete metric---basis points of implied policy rate difference---that translates statistical accuracy into a policymaker-relevant quantity.

  \item \textbf{Cross-country evidence.}  Results for four Euro-area economies establish external validity and suggest that a centralised ECB nowcasting system is feasible.

  \item \textbf{Reproducibility.}  All data come from public APIs, all code is open-source, and the expanding-window evaluation is fully replicable.
\end{enumerate}


\section*{Policy Recommendations}

Based on our findings, we offer three recommendations to statistical agencies and central banks:

\begin{enumerate}
  \item \textbf{Adopt Ridge regression for operational BoP nowcasting.}  Ridge is simple, fast, interpretable, and outperforms all competitors in our evaluation.  It can be deployed with minimal infrastructure and updated automatically as new data arrive.

  \item \textbf{Incorporate text/sentiment data.}  BdF-type business surveys and NLP-based central bank sentiment scores provide high signal at low cost.  These should be prioritised over satellite data, which are more expensive and contribute less.

  \item \textbf{Invest in real-time data infrastructure.}  The policy value of better nowcasts depends on their availability: a perfect estimate published with the same delay as the official data has zero marginal value.  Agencies should develop automated pipelines that update nowcasts daily as new indicators are released.
\end{enumerate}


\section*{Limitations}

Several limitations should be acknowledged:

\begin{itemize}
  \item \textbf{ECB text coverage.}  Our FinBERT-based sentiment scores rely on scraped ECB press conference transcripts, of which we successfully retrieved only~13\%.  Direct access to ECB communication archives or use of the ECB's machine-readable data feeds would substantially improve this pipeline.

  \item \textbf{Satellite data.}  Nighttime light intensity, our sole satellite proxy, is partially synthetic (pre-2012 data are interpolated from IP and electricity consumption).  AIS vessel tracking data or port-level container counts would provide more granular physical-activity measures.

  \item \textbf{LSTM implementation.}  The LSTM's failure may partly reflect our implementation choices (two-layer architecture, 32 hidden units).  More sophisticated architectures---Temporal Fusion Transformers \citep{lim2021temporal}, GRU variants, or pre-trained models---might perform better, though sample-size constraints remain binding.

  \item \textbf{Single-country focus.}  While Chapter~3 provides cross-country evidence, the ablation study (Chapter~2) is France-specific.  Extending it to other countries would strengthen the external validity of the alternative data findings.

  \item \textbf{Financial account.}  We focus on the current account (goods trade balance); the BoP's financial account---which records portfolio flows, FDI, and derivatives---poses different challenges and may require different modelling approaches.
\end{itemize}


\section*{Future Research}

Five directions for future work emerge:

\begin{enumerate}
  \item \textbf{Transformer-based architectures.}  Pre-trained language models finetuned on economic text \citep{devlin2019bert} could provide both feature extraction (replacing FinBERT) and time-series forecasting capabilities.  The recently proposed Chronos models for time-series foundation modelling are a promising direction.

  \item \textbf{Real-time data vintage.}  Our pseudo-real-time evaluation uses final-vintage data.  Using actual real-time vintages (with revisions) from the ECB's Real-Time Database would provide a more stringent test.

  \item \textbf{Forecast combination.}  Combining Ridge with Bridge or DFM forecasts (which exploit different information structures) may outperform any individual method.  Optimal combination weights could be estimated following \citet{stock2004combination}.

  \item \textbf{Financial account nowcasting.}  Portfolio flows, FDI, and TARGET2 balances are even less timely than current account data and are potentially more amenable to alternative data approaches (e.g., using EPFR fund flow data or SWIFT transaction records).

  \item \textbf{Causal identification.}  Our Taylor rule exercise is correlational.  A structural VAR or narrative identification strategy could establish whether better BoP information \emph{causally} affects policy decisions, rather than merely being correlated with them.
\end{enumerate}
